% ============================================
% Supplementary Note 3: Robustness of the Simulated Dominance Transition to Model Assumptions
% ============================================
\subsection*{Supplementary Note 3: Robustness of the simulated Dominance transition to model assumptions}
\label{sec:suppl:note3}

\rev{The baseline gLV model and the interaction-strength parameter $\mu$ are defined in Supplementary Methods. Here, $\mu$ provides a minimal one-parameter control of effective interaction strength that jointly sets the average inhibitory effect and the heterogeneity of pairwise coefficients~\citep{Hu2022,Marsland2019,Ratzke2020,Hu2025,Allesina2012,Grilli2017}. We test whether the simulated transition from Mixture-dominated to Dominance/Restructuring outcomes with increasing $\mu$ is robust to the assumptions of that baseline, including the coefficient distribution, demographic heterogeneity in growth rates and carrying capacities, community size, interaction sign structure, and the decomposition of $\mu$ into coefficient mean and variance.}

\subsubsection*{\rev{Baseline gLV sensitivity analyses}}

\rev{We first tested whether heterogeneity in intrinsic growth rates or carrying capacities altered the qualitative transition. Allowing species-level variation in growth rates or carrying capacities changed the detailed outcome frequencies but did not remove the increase in Dominance with the interaction-strength parameter $\mu$ (Supplementary Figs.~4--5).}

\rev{We next tested whether the qualitative transition depended on the shape of the coefficient distribution. We compared three distributions with matched coefficient mean $\mu$: uniform ranging from 0 to $2\mu$ (primary analysis), Gaussian with mean $\mu$ and standard deviation $\mu/\sqrt{3}$, and Gamma with mean $\mu$ and shape parameter 3. All three distributions were parameterized to have the same mean $\mu$ and matched variance $\mu^2/3$ (for the uniform distribution $U[0,2\mu]$, $\text{var} = (2\mu)^2/12 = \mu^2/3$; for the Gaussian, exact; for the Gamma, exact). The qualitative transition from Mixture-dominated to Dominance/Restructuring outcomes with increasing interaction-strength parameter $\mu$ persisted across all tested distributions (Supplementary Fig.~6).}

We also tested whether community size (number of species per parental community) affects the frequency of Dominance outcomes. Simulations were performed with parental community sizes ranging from 4 to 48 species per community (\rev{Extended Data Fig.~5}). In this ablation, the total simulated species pool scaled with parental-community richness as four non-overlapping parental communities ($N=16, 24, 36, 48, 96,$ or $192$ species for 4, 6, 9, 12, 24, or 48 species per community, respectively). Dominance frequency changed little across community sizes ($\sim$14--26\% at $\mu=0.3$, $\sim$56--67\% at $\mu=0.6$, $\sim$69--78\% at $\mu=0.8$), while Mixture decreased and Restructuring increased with larger species pools. This pattern reflects increased opportunity for competitive exclusion cascades in larger communities, but Dominance remained common at moderate-to-high interaction-strength parameter values.

\subsubsection*{\rev{Beneficial-interaction and reciprocal pair-coupling extensions}}

\rev{Because our baseline gLV ensemble is restricted to non-negative, growth-inhibiting coefficients, we also tested additional interaction structures while applying the same assembly, coalescence, and outcome-classification pipeline. First, we introduced a mixed-sign facilitative-tail model in which off-diagonal coefficients were sampled from $\alpha_{ij}\sim U[-f\mu,(2+f)\mu]$, preserving $E[\alpha_{ij}]=\mu$ while allowing a small fraction of coefficients to be negative under our sign convention (positive ecological interactions). To prevent runaway growth in the presence of facilitation, we added density-dependent self-limitation, $\dot n_i=n_i(1-(A n)_i-0.1n_i^2)$. Across $f \in \{0,0.10,0.20,0.40,0.80\}$ and $\mu \in \{0.30,0.60,0.80\}$, Dominance still increased with $\mu$, although increasing the facilitative tail shifted some low-$\mu$ outcomes from Mixture toward Dominance and Restructuring (Supplementary Fig.~39).}

\rev{Second, we varied the reciprocal pair-coupling structure between coefficients $\alpha_{ij}$ and $\alpha_{ji}$. For a fraction $|p|$ of unordered species pairs, coefficients were constrained to follow an antisymmetric exploitation structure when $p<0$ ($\alpha_{ji}=-\alpha_{ij}$) or a symmetric competition structure when $p>0$ ($\alpha_{ji}=\alpha_{ij}$). Remaining pairs were drawn iid from the baseline competitive distribution. This pair-coupling sweep tests whether the Dominance increase depends on independent directed interaction coefficients. The qualitative $\mu$-dependent increase in Dominance persisted across the reciprocal pair-coupling sweep and in a finer sweep from $p=-1$ to $p=+1$ (Supplementary Fig.~40).}

\rev{Third, to address the specific mutualism question, we converted a controlled fraction of unordered species pairs into weak bidirectional mutualisms. For $q \in \{0,0.10,0.20,0.30,0.40\}$ of species pairs, both reciprocal coefficients were sampled from $\alpha_{ij},\alpha_{ji}\sim U[-0.2\mu,0]$, while all remaining off-diagonal coefficients were sampled from the baseline competitive distribution $U[0,2\mu]$. We used the same density-dependent self-limitation and the same assembly, coalescence, and classification pipeline as in the mixed-sign facilitative-tail analysis. Weak mutualistic pairs did not eliminate the interaction-strength-dependent Dominance trend: at $\mu=0.60$, Dominance was 58.5\%, 70.2\%, 69.8\%, 67.3\%, and 58.7\% for $q=0,0.10,0.20,0.30,0.40$, respectively, and at $\mu=0.80$ it remained 69.4--77.5\% across the same sweep (Supplementary Fig.~41).}

\subsubsection*{\rev{Mean and heterogeneity of interaction coefficients}}

\rev{Fourth, to test which components of this one-parameter interaction-strength axis are required, we separated the mean and variance of the interaction-coefficient distribution. Off-diagonal coefficients were sampled from $\alpha_{ij}\sim U[m-h,m+h]$ for $m \in \{0,0.2,0.4,0.6,0.8\}$ and $h \in \{0,0.2,0.4,0.6,0.8\}$, giving $\mathrm{std}(\alpha_{ij})=h/\sqrt{3}$. Negative coefficients, corresponding to positive ecological interactions under our sign convention, occur when $h>m$. When $h=0$, all outcomes were Mixture even at $m=0.8$, showing that coefficient mean alone is not sufficient to generate Dominance. Increasing $h$ at fixed $m$ increased Dominance and same-parent fate concordance. The strongest Dominance values occurred when both the coefficient mean and coefficient variance were high. Across the 25 cells, Dominance correlated more strongly with $h$ (and therefore coefficient standard deviation) than with $m$ (descriptive Pearson $r=0.78$ versus $0.53$; Supplementary Fig.~42). Thus, both average inhibitory effect and coefficient heterogeneity contribute to the effective interaction-strength axis captured by the baseline $\mu$ parameter. These sensitivity checks indicate that the main simulation result is not restricted to iid non-negative coefficients. Strongly mutualistic or facilitation-dominated systems remain outside the baseline model scope~\citep{Tikhonov2016,Lechon2021}.}
